% =============================================================================
% Chapter 14: ML Operations — ML Systems Lecture Slides
% =============================================================================
% IMAGES NEEDED (SVGs converted to PDF):
%   - cicd-pipeline-ml.pdf       - monitoring-dashboard-ml.pdf
%   - drift-detection-flow.pdf   - canary-deployment.pdf
%   - maturity-model.pdf         - degradation-equation.pdf
% =============================================================================
\documentclass[aspectratio=169, 12pt]{beamer}
\usepackage{../../assets/beamerthememlsys}

\mlsyssetup{
  volume   = {Volume I},
  chapter  = {Chapter 14},
  logo = {../../assets/img/logo-mlsysbook.png},
  instlogo = {../../assets/img/logo-harvard.png},
  chaptertitle = {ML Operations},
}

% --- Fonts ---
\usepackage{fontspec}
\setsansfont{Helvetica Neue}[
  BoldFont={Helvetica Neue Bold},
  ItalicFont={Helvetica Neue Italic},
  BoldItalicFont={Helvetica Neue Bold Italic},
]
\setmonofont{JetBrains Mono}[Scale=0.85]

% --- Packages ---
\usepackage{booktabs}
\usepackage{amsmath}

% --- Image paths ---
\graphicspath{
  {images/}
}

% --- Chapter-specific macros ---
\newcommand{\DAM}{D\raisebox{0.08em}{\tiny$\bullet$}A\raisebox{0.08em}{\tiny$\bullet$}M}

% --- Helper: safe image include ---
\newcommand{\safeimg}[2][width=\textwidth,keepaspectratio]{%
  \IfFileExists{images/#2}{\includegraphics[#1]{#2}}{%
    \IfFileExists{#2}{\includegraphics[#1]{#2}}{%
      \fbox{\parbox[c][2.5cm][c]{0.85\linewidth}{\centering\footnotesize\textcolor{midgray}{[Missing image]}}}%
    }%
  }%
}

% --- Section count for navigation (must match actual \section{} count) ---
\setcounter{mlsystotalsections}{8}

\title{ML Operations}
\author{Vijay Janapa Reddi}
\institute{Harvard University}
\date{}

\begin{document}

% =============================================================================
% TITLE SLIDE
% =============================================================================
\mlsystitle{ML Operations}{Keeping Models Correct in Production}{cover_ml_ops.png}

% =============================================================================
% LEARNING OBJECTIVES
% =============================================================================
\begin{frame}{Learning Objectives}
\note{
% -- LINK: Connects to deployment from Ch13
Students deployed models in the previous chapter. This slide reframes
the question: deployment is the beginning of operations, not the end.

% -- NARRATE: Walk through objectives with production framing
Read each objective aloud. Pause after ``silent failures'': ``This is the
central challenge --- your model is wrong and nobody knows.''
ANALOGY: Traditional software crashes like a car engine seizing. ML
systems degrade like carbon monoxide --- invisible, gradual, deadly.

% -- ENGAGE: Surface students' production experience
Ask: ``How many of you have deployed a model and checked on it a month
later?'' Follow up: ``What did you find?'' [Expected: most never checked]

% -- WARN: Students assume deployment is the finish line
Common error: students treat deployment as the final step.
Correct framing: deployment is where the hardest engineering begins ---
data changes, users change, the world changes.

% -- FLEX: [CORE] Sets the agenda for the entire lecture
[CORE] Do not skip --- these objectives structure every section.
IF AHEAD: Ask students to rank which objective seems hardest.
IF SHORT: Read objectives without discussion, move to next slide.
}

\small
\begin{enumerate}
  \item Explain why \textbf{silent failures} distinguish ML from traditional software
  \item Analyze \textbf{technical debt} patterns: boundary erosion, correction cascades, data dependencies
  \item Design \textbf{feature stores} and \textbf{CI/CD pipelines} for training-serving consistency
  \item Apply the \textbf{retraining staleness model} to optimize retraining intervals
  \item Quantify \textbf{data drift} using \textbf{PSI} and \textbf{KL divergence}
  \item Compare \textbf{canary}, \textbf{blue-green}, and \textbf{shadow} deployment strategies
  \item Implement structured \textbf{incident response} for ML-specific failure modes
  \item Evaluate organizational \textbf{MLOps maturity levels}
\end{enumerate}

\end{frame}

\begin{frame}{Visual Language}
\note{
% -- NARRATE: Orient students to the color system for this deck
Point to each card: ``Blue for compute, green for data flow, orange for
routing and scheduling, red for errors and bottlenecks. In this chapter,
you will see a lot of orange (scheduling retraining) and red (drift and
silent failures).''

% -- FLEX: [OPTIONAL] Returning students already know this
[OPTIONAL] If students have seen previous decks, spend 30 seconds max.
IF SHORT: Skip entirely --- students already internalized the palette.
}

\small
Throughout this course, colors carry meaning:

\vspace{0.3cm}
\begin{columns}[T]
  \begin{column}{0.45\textwidth}
    \begin{mlsyscard}{computestroke}
    \textbf{Blue} --- Compute / Processing\\
    {\footnotesize GPU ops, forward/backward pass, inference}
    \end{mlsyscard}
    \vspace{0.15cm}
    \begin{mlsyscard}{datastroke}
    \textbf{Green} --- Data / Memory\\
    {\footnotesize Data flow, caches, healthy paths}
    \end{mlsyscard}
  \end{column}
  \begin{column}{0.45\textwidth}
    \begin{mlsyscard}{routingstroke}
    \textbf{Orange} --- Routing / Scheduling\\
    {\footnotesize Load balancers, batch windows}
    \end{mlsyscard}
    \vspace{0.15cm}
    \begin{mlsyscard}{errorstroke}
    \textbf{Red} --- Error / Cost / Bottleneck\\
    {\footnotesize Loss, decode phase, waste}
    \end{mlsyscard}
  \end{column}
\end{columns}
\end{frame}


% =============================================================================
\section{The Operational Mismatch}
% =============================================================================

\begin{frame}{The Operational Mismatch}
\note{
% -- LINK: From learning objectives to the first concrete failure
The objectives promised ``silent failures.'' This slide delivers the first
one: a real production system degrading while every dashboard stayed green.

% -- NARRATE: Tell the ridesharing story as a narrative
``A ridesharing company deploys a demand model. 94\% accuracy at launch.
Week 4: 88\%. Week 8: drivers dispatched to wrong areas. The entire time,
every infrastructure metric --- CPU, memory, uptime --- was green.''
Point to the DevOps vs MLOps table: ``DevOps monitors crashes. MLOps
monitors drift. Different disease, different diagnostics.''

% -- ENGAGE: Force students to think about detection
Ask: ``What metrics would have caught the 6-point accuracy drop at week 2
instead of week 8?'' [Expected: input distribution stats, not just uptime]

% -- WARN: Students conflate infrastructure health with model health
Common error: students assume green dashboards mean the system is working.
Correct framing: infrastructure health and model health are orthogonal ---
a perfectly available system can be perfectly wrong.

% -- FLEX: [CORE] This example motivates the entire chapter
[CORE] The ridesharing story is the anchor for every subsequent concept.
IF AHEAD: ``What was the root cause?'' (Competitor promotion shifted demand.)
IF SHORT: Skip the table, keep the narrative and the punchline.
}

\small
\begin{columns}[T]
  \begin{column}{0.55\textwidth}
    A ridesharing demand model: 94\% accuracy at deploy.

    \vspace{0.15cm}
    \begin{itemize}\setlength\itemsep{1pt}
      \item \textbf{Week 1}: Excellent performance
      \item \textbf{Week 4}: Accuracy drops to 88\%
      \item \textbf{Week 8}: Driver dispatch is inefficient
      \item \textbf{All infrastructure metrics}: \textcolor{datastroke}{\textbf{Green}}
    \end{itemize}

    \vspace{0.15cm}
    \begin{mlsyscard}{errorstroke}
    A competitor launched a promotion that shifted user behavior.
    The model needed retraining \textbf{6 weeks ago}.
    \end{mlsyscard}
  \end{column}
  \begin{column}{0.42\textwidth}
    \renewcommand{\arraystretch}{1.2}
    {\footnotesize
    \begin{tabular}{@{}lll@{}}
      \toprule
      & \textbf{DevOps} & \textbf{MLOps} \\
      \midrule
      Monitors & Uptime & Accuracy \\
      Failure & Crash & Drift \\
      Signal  & Error log & Statistics \\
      Fix     & Patch code & Retrain \\
      \bottomrule
    \end{tabular}
    }
  \end{column}
\end{columns}

\end{frame}

\begin{frame}{Why MLOps Exists}
\note{
% -- LINK: From the ridesharing failure to the formal definition
The ridesharing model failed because nobody closed the feedback loop
between production data and model behavior. MLOps is that feedback loop.

% -- NARRATE: Read the definition, then unpack the three interfaces
Read the crimson card aloud, emphasizing ``measurable production drift.''
Then point to each row: ``Data-Model: are training and serving seeing
the same features? Model-Infra: is the environment identical?
Prod-Monitor: when do we retrain?''

% -- ENGAGE: Test the definition
Ask: ``Is MLOps just DevOps plus a model registry?'' [Expected: some say
yes.] ``What does DevOps lack?'' [Answer: data validation, drift detection,
retraining triggers --- the three interfaces in this table.]

% -- WARN: Students reduce MLOps to tooling
Common error: students think MLOps = installing MLflow or Kubeflow.
Correct framing: MLOps is a discipline for managing three interfaces;
tools implement it, but the discipline is about closing feedback loops.

% -- FLEX: [CORE] The three-interface framework recurs throughout
[CORE] Students will reference Data-Model, Model-Infra, and Prod-Monitor
in every subsequent section.
IF AHEAD: ``Which interface failed in the ridesharing example?'' (Prod-Monitor.)
IF SHORT: Skip the engage question, keep the definition and table.
}

\small
\begin{mlsyscard}{crimson}
\textbf{MLOps:} The engineering discipline that closes the feedback loop between
model behavior and data reality by automating retraining, validation, and
deployment in response to measurable production drift.
\end{mlsyscard}

\vspace{0.15cm}
\textbf{Three Critical Interfaces:}

\vspace{0.1cm}
{\footnotesize
\renewcommand{\arraystretch}{1.15}
\begin{tabular}{@{}lll@{}}
  \toprule
  \textbf{Interface} & \textbf{Challenge} & \textbf{Solution} \\
  \midrule
  \textcolor{datastroke}{\textbf{Data-Model}}       & Feature consistency     & Feature stores \\
  \textcolor{computestroke}{\textbf{Model-Infra}}   & Environment parity      & Registries + containers \\
  \textcolor{routingstroke}{\textbf{Prod-Monitor}}   & Retraining cadence      & Drift detection \\
  \bottomrule
\end{tabular}
}

\end{frame}

% --- ACTIVE LEARNING 1: Predict ---
\begin{frame}{Predict: What Breaks First?}
\note{
% -- LINK: From the operational mismatch to predicting failure modes
Students just saw one failure (competitor promotion). Now they must
generate their own failure taxonomy before the chapter reveals it.

% -- NARRATE: Frame the exercise with urgency
``No code changed. Six months passed. What broke?'' Pause for effect.

% -- ENGAGE: Think-Write-Share with 60-second timer
Students write 3 failure modes individually (60 seconds), then compare
with a neighbor. Cold-call 2--3 students. Do NOT reveal answers yet ---
the technical debt section will formalize their intuitions.
[Expected answers: data drift, feature staleness, upstream schema change]

% -- FLEX: [CORE] Primes the entire technical debt section
[CORE] This prediction activates prior knowledge and creates curiosity.
IF SHORT: Reduce to 30-second think, skip the neighbor comparison.
}

\centering
\vspace{0.8cm}
{\Large\bfseries Think--Write--Share}

\vspace{0.4cm}
{\large You deploy a model with 94\% accuracy.\\
Six months later, no code has changed.\\[0.2cm]
\alert{What are the three most likely failure modes?}}

\vspace{0.4cm}
{\normalsize Write down 3 failure modes. \textcolor{midgray}{(60 seconds)}}

\vspace{0.3cm}
{\small\textcolor{lightgray}{Turn to a neighbor and compare.}}

\end{frame}

% =============================================================================
\section{Technical Debt}
% =============================================================================

\begin{frame}{The 5\% Problem in Production}
\note{
% -- LINK: Students predicted failure modes; now we formalize them
The failure modes students generated map directly onto Sculley et al.'s
six debt patterns.

% -- NARRATE: The 5\% framing is the hook
``ML code is 5\% of the system. The other 95\% --- data validation,
feature engineering, configuration, monitoring --- is where technical
debt accumulates.'' Point to the CACHE principle: ``Change Anything
Changes Everything.'' Point to the break-even card: ``80 hours to
automate vs 4 hrs/week manual. Break-even at 20 weeks.''

% -- ENGAGE: Connect back to the ridesharing failure
Ask: ``Which of these six debt patterns caused the ridesharing failure?''
[Expected: Data debt (quality drift) + Feedback loops (stale retraining)]

% -- WARN: Students underestimate data debt relative to code debt
Common error: students focus on code quality and ignore data quality.
In ML systems, data debt accumulates faster and is harder to detect
because data has no compiler.

% -- FLEX: [CORE] The six patterns are referenced throughout the chapter
[CORE] CACHE principle and break-even calculation recur in later sections.
IF AHEAD: ``Which pattern is hardest to detect? Why?'' (Feedback loops.)
IF SHORT: Skip the break-even card, keep CACHE and the six patterns.
}

\small
\begin{columns}[T]
  \begin{column}{0.48\textwidth}
    \textbf{Six ML Debt Patterns:}

    \vspace{0.15cm}
    {\footnotesize
    \renewcommand{\arraystretch}{1.1}
    \begin{tabular}{@{}ll@{}}
      \toprule
      \textbf{Pattern} & \textbf{Mechanism} \\
      \midrule
      Boundary erosion & CACHE principle \\
      Correction cascades & Fix $\to$ break \\
      Feedback loops & Self-reinforcing \\
      Data debt & Quality drift \\
      Config debt & Parameter sprawl \\
      Pipeline debt & Fragile workflows \\
      \bottomrule
    \end{tabular}
    }

    \vspace{0.15cm}
    \mlsysalert{CACHE:}{Change Anything Changes Everything.}
  \end{column}
  \begin{column}{0.48\textwidth}
    \begin{mlsyscard}{errorstroke}
    {\footnotesize
    Manual: 4 hrs/week\\
    Pipeline: 80 hrs (one-time)\\
    Break-even: \textbf{20 weeks}\\
    After 1 yr: \textbf{100\%} maintenance
    }
    \end{mlsyscard}

    \vspace{0.05cm}
    {\footnotesize Automation is about \textbf{capacity}.}
  \end{column}
\end{columns}
\mlsyscite{Sculley et al., NeurIPS 2015}

\end{frame}

\begin{frame}{Real-World Debt: Case Studies}
\note{
% -- LINK: From abstract debt patterns to real-world consequences
The previous slide named six patterns. This slide shows four companies
that paid the price.

% -- NARRATE: Lead with Zillow, the most dramatic
``Zillow's algorithm overvalued homes, then used those purchases as future
training data --- a textbook correction cascade. Cost: over \$500 million.''
Then YouTube: ``Recommendations shaped clicks, clicks became training data,
training data reinforced recommendations --- a feedback loop.''

% -- ENGAGE: Pattern matching exercise
Ask: ``Which debt pattern does each case represent?''
[Zillow = correction cascade, YouTube = feedback loop, Tesla = undeclared
consumer, Facebook = config debt]

% -- WARN: Students think these are edge cases
Common error: students assume these failures happen only at Big Tech scale.
Correct framing: the patterns are scale-invariant --- a 3-person startup
can have correction cascades in a Jupyter notebook pipeline.

% -- FLEX: [OPTIONAL] Case studies reinforce but are not structurally required
[OPTIONAL] These are motivational, not foundational.
IF SHORT: Cover only Zillow and YouTube, skip Tesla and Facebook.
}

\scriptsize
\renewcommand{\arraystretch}{1.0}
\begin{tabular}{@{}llp{4.0cm}l@{}}
  \toprule
  \textbf{Company} & \textbf{Debt Type} & \textbf{What Happened} & \textbf{Cost} \\
  \midrule
  \textbf{Zillow} & Correction cascade & Overvalued homes; errors fed training & \textcolor{errorstroke}{\$500M+} \\
  \textbf{YouTube} & Feedback loop & Recs shaped clicks $\to$ training data & Content harm \\
  \textbf{Tesla} & Undeclared consumer & OTA updates changed subsystems & Safety risk \\
  \textbf{Facebook} & Config debt & Opaque settings drove ranking & Misalignment \\
  \bottomrule
\end{tabular}

\vspace{0.1cm}
\small
\mlsysalert{Lesson:}{Predictable consequences of probabilistic systems without operational infrastructure.}

\end{frame}

\begin{frame}{Quick Check}
\note{
% -- NARRATE: Cue the retrieval exercise
``Close your notes. 30 seconds. What are the three critical interfaces?''
Pause silently. After 30 seconds: ``Data-Model, Model-Infra, Prod-Monitor.''

% -- FLEX: [CORE] Distributed retrieval cements the three-interface model
[CORE] Reinforces the framework before Dev Infrastructure builds on it.
IF SHORT: Reduce to 15 seconds, give the answer verbally.
}

\centering
\Large\bfseries Quick check --- no peeking.\\[0.5cm]
\normalsize What are the three critical interfaces that MLOps must manage?\\[0.3cm]
{\small 30 seconds --- then we continue.}
\end{frame}



% =============================================================================
\section{Dev Infrastructure}
% =============================================================================

\begin{frame}{Five Foundational Principles}
\note{
% -- LINK: From technical debt to the principles that prevent it
Technical debt accumulates when principles are missing. These five
principles are the antidote --- each maps to infrastructure.

% -- NARRATE: Walk through the table row by row
``Reproducibility: version everything --- code, data, config, environment.''
Point to the equation: ``Model Output = f(Code\_v, Data\_v, Config\_v,
Environment\_v). Missing any version makes the model unreproducible.''

% -- ENGAGE: Test the reproducibility equation
Ask: ``Which of these four versions is most commonly forgotten?''
[Expected: Data\_v or Environment\_v --- students rarely version datasets]

% -- WARN: Students think git covers reproducibility
Common error: students version code and assume the model is reproducible.
Correct framing: identical code on different data or library versions
produces different models. All four inputs must be versioned.

% -- FLEX: [CORE] These five principles are the evaluation framework
[CORE] Every subsequent infrastructure slide maps back to one of these.
IF AHEAD: ``Which principle does a feature store implement?'' (Consistency.)
IF SHORT: Read the table without row-by-row narration.
}

\footnotesize
\renewcommand{\arraystretch}{1.2}
\begin{tabular}{@{}lp{4.2cm}l@{}}
  \toprule
  \textbf{Principle} & \textbf{Core Insight} & \textbf{Key Metric} \\
  \midrule
  \textbf{1. Reproducibility}     & Version all artifacts (code, data, config, env) & Artifact hash \\
  \textbf{2. Separation}          & Independent layer evolution & Coupling score \\
  \textbf{3. Consistency}         & Training $=$ Serving & Feature skew rate \\
  \textbf{4. Observable Degrad.}  & Make silent failures visible & Time to detection \\
  \textbf{5. Cost-Aware Auto.}    & Optimize total cost & Net retrain value \\
  \bottomrule
\end{tabular}

\vspace{0.2cm}
\mlsysconcept{Model Output =}{$f(\text{Code}_v, \text{Data}_v, \text{Config}_v, \text{Environment}_v)$}

\vspace{0.1cm}
{\footnotesize Missing \textbf{any} version makes the model \textbf{unreproducible}.}

\end{frame}

\begin{frame}{Feature Stores: Eliminating Training-Serving Skew}
\note{
% -- LINK: From Principle 3 (Consistency) to its implementation
The five principles named Consistency as training = serving. The feature
store is the infrastructure that enforces it.

% -- NARRATE: The session\_length story makes skew visceral
``A data scientist computes session\_length using wall-clock time in Python.
An engineer reimplements it in Java using processing time. 8\% accuracy
drop nobody can explain.'' Point to the Skew Cost card: ``1\% skew on 1M
queries/day at \$0.10 each = \$365,000/year in silent losses.''

% -- ENGAGE: Detection challenge
Ask: ``How would you detect this skew without a feature store?''
[Expected: compare training vs serving distributions --- expensive and
manual. The feature store eliminates the problem at the source.]

% -- WARN: Students think schema checks are sufficient
Common error: students assume matching column names means features are
consistent. Correct framing: skew lives in the computation, not the
schema --- same column name, different computation logic.

% -- FLEX: [CORE] Feature stores are central to the Data-Model interface
[CORE] The \$365K/year number recurs in Key Takeaways.
IF AHEAD: ``What about feature freshness --- how stale can features be?''
IF SHORT: Skip the skew type table, keep the story and the cost card.
}

\small
\begin{columns}[T]
  \begin{column}{0.55\textwidth}
    \textbf{The Problem:}\\
    {\footnotesize Training computes \texttt{session\_length} with wall-clock time.\\
    Serving uses processing time. 8\% accuracy drop.}

    \vspace{0.15cm}
    \textbf{The Solution:}\\
    {\footnotesize Compute features \textbf{once}, serve to both pipelines.}

    \vspace{0.15cm}
    \begin{mlsyscard}{datastroke}
    {\footnotesize
    \textbf{Skew Cost} $=$ Error Rate $\times$ Query Vol $\times$ Impact\\[0.05cm]
    1\% skew on 1M queries/day $\times$ \$0.10 each\\
    $=$ \textbf{\$365,000/year} in silent losses
    }
    \end{mlsyscard}
  \end{column}
  \begin{column}{0.42\textwidth}
    {\footnotesize
    \renewcommand{\arraystretch}{1.1}
    \begin{tabular}{@{}ll@{}}
      \toprule
      \textbf{Skew Type} & \textbf{Detection} \\
      \midrule
      Preprocessing & Stats comparison \\
      Missing data  & Schema validation \\
      Time features & Timestamp checks \\
      Library drift & Env hash compare \\
      \bottomrule
    \end{tabular}
    }

    \vspace{0.15cm}
    \mlsysinsight{Result:}{Feature stores eliminate 5--15\% accuracy loss.}
  \end{column}
\end{columns}

\end{frame}

\begin{frame}{ML CI/CD Pipelines}
\note{
% -- LINK: From feature stores (one component) to the full pipeline
Feature stores solve consistency. CI/CD pipelines orchestrate the entire
flow: data validation, training, evaluation, deploy.

% -- NARRATE: Walk through the three phases in the diagram
Point to each phase: ``Data validation takes 20--30\% of pipeline time.
Training is 40--60\%. Then canary deployment with graduated rollout.''
Emphasize the fail loop: ``If model evaluation fails, reject and go back
to data validation. This loop is the pipeline's immune system.''

% -- ENGAGE: Compare pipeline durations
Ask: ``How long does a traditional CI/CD pipeline take? How long for ML?''
[Expected: traditional 2--5 min, ML 15--45 min. The 3--9x ratio surprises.]

% -- WARN: Students skip data validation to speed up pipelines
Common error: students cut data validation to reduce pipeline time,
then spend weeks debugging silent accuracy drops.

% -- FLEX: [CORE] The pipeline structure frames deployment strategies later
[CORE] Canary deployment in this diagram connects to the Deployment section.
IF AHEAD: ``What happens if data validation has a false negative?''
IF SHORT: Show the image, narrate the three phases briefly.
}

% --- Layout: FULL-WIDTH IMAGE ---
\centering
\safeimg[width=0.92\textwidth,height=4.5cm,keepaspectratio]{cicd-pipeline-ml.pdf}

\vspace{0.1cm}
\small
\begin{columns}[T]
  \begin{column}{0.30\textwidth}
    \centering
    \textcolor{datastroke}{\textbf{Data Validation}}\\
    {\scriptsize 20--30\% of pipeline time}
  \end{column}
  \begin{column}{0.30\textwidth}
    \centering
    \textcolor{computestroke}{\textbf{Model Training}}\\
    {\scriptsize 40--60\% of pipeline time}
  \end{column}
  \begin{column}{0.30\textwidth}
    \centering
    \textcolor{routingstroke}{\textbf{Canary Deploy}}\\
    {\scriptsize Graduated rollout}
  \end{column}
\end{columns}

\end{frame}

% =============================================================================
\section{Drift \& Monitoring}
% =============================================================================

\mlsysfocus{The Degradation Equation}{%
$\text{Acc}(t) \approx \text{Acc}_0 - \lambda \cdot D(P_t \| P_0)$%
}

\begin{frame}{The Rotting Asset Curve}
\note{
% -- LINK: From the Degradation Equation to its visual consequences
The equation was just shown on the focus slide. This chart shows what
it looks like in practice over time.

% -- NARRATE: Walk through both curves in the diagram
Point to the orange sawtooth: ``Scheduled retraining every 90 days. See
how accuracy dips below the SLA threshold before each retrain? Revenue
lost.'' Point to the green curve: ``Trigger-based retraining: when drift
exceeds a threshold, retrain immediately. Never violates SLA.''

% -- ENGAGE: Domain-specific strategy selection
Ask: ``Which strategy for a fraud detection model with 10M daily
transactions?'' [Expected: trigger-based, because each day below SLA
costs real money at that query volume.]

% -- WARN: Students default to scheduled because it is simpler
Common error: students prefer fixed-interval retraining because it is
predictable. Correct framing: drift is not periodic. A competitor launch
or holiday season can shift distributions at any time.

% -- FLEX: [CORE] The sawtooth vs trigger visual anchors the retraining exercise
[CORE] The next active learning slide uses the $T^*$ formula from this.
IF AHEAD: ``What if retraining is cheap? Does scheduled ever win?''
IF SHORT: Show both curves, state the punchline, skip the question.
}

% --- Layout: FULL-WIDTH IMAGE ---
\centering
\safeimg[width=0.92\textwidth,height=4.8cm,keepaspectratio]{degradation-equation.pdf}

\vspace{0.1cm}
\footnotesize
\begin{columns}[T]
  \begin{column}{0.45\textwidth}
    \textcolor{routingstroke}{\textbf{Scheduled}}: Fixed interval, may violate SLA
  \end{column}
  \begin{column}{0.45\textwidth}
    \textcolor{datastroke}{\textbf{Trigger-based}}: Adaptive, higher avg accuracy
  \end{column}
\end{columns}

\end{frame}

\begin{frame}{Data Drift vs.\ Concept Drift}
\note{
% -- LINK: From the Degradation Equation to diagnosing the cause
The equation tells us accuracy is dropping. This slide asks: why?
Data drift and concept drift are two different diseases.

% -- NARRATE: Distinguish the two types with concrete examples
``Data drift: the input distribution P(X) shifts. Your fraud model was
trained on pre-holiday transactions; now it sees Black Friday patterns.
Retrain on fresh data.'' ``Concept drift: the relationship P(Y|X) changes.
COVID changed what 'normal health metrics' meant --- same inputs, different
correct labels. You need new labels, not just new data.''

% -- ENGAGE: Classification exercise
Ask: ``A competitor launches a new product. Data drift or concept drift?''
[Expected: data drift --- user behavior shifts, but the relationship
between features and outcomes is unchanged.]

% -- WARN: Students retrain without diagnosing the cause
Common error: students detect drift and immediately retrain on latest data.
If the drift is concept drift, retraining on shifted data entrenches the
shift --- the model learns the wrong labels.

% -- FLEX: [CORE] The data/concept distinction determines the correct response
[CORE] This distinction drives monitoring strategy in the next slides.
IF AHEAD: ``Can both types occur simultaneously?''
IF SHORT: Skip the engage question, keep the definitions and alert box.
}

% --- Layout: FULL-WIDTH IMAGE ---
\centering
\safeimg[width=0.88\textwidth,height=4.5cm,keepaspectratio]{drift-detection-flow.pdf}

\vspace{0.15cm}
\mlsysalert{Key:}{Diagnose \emph{which} distribution changed before retraining. Retraining on shifted data entrenches the shift.}

\end{frame}

\begin{frame}{Quantifying Drift: PSI}
\note{
% -- LINK: From diagnosing drift type to measuring drift magnitude
We know data drift and concept drift exist. PSI answers: how much?

% -- NARRATE: Walk through the formula term by term
``PSI sums over bins. For each bin: take the difference between actual
and expected proportions, multiply by the log ratio. A 3-percentage-point
shift yields PSI = 0.029 --- below the 0.1 warning threshold.''
Point to the threshold table: ``PSI for categorical, KS for continuous,
KL for full distributions.''

% -- ENGAGE: Threshold interpretation
Ask: ``Your model's PSI is 0.15. What do you do?''
[Expected: investigate which features shifted --- do not immediately
retrain. 0.15 is between warning (0.1) and significant (0.2).]

% -- WARN: Students monitor outputs instead of inputs
Common error: students track output accuracy and wait for it to drop.
Correct framing: input drift precedes output degradation by the length
of the ground-truth feedback loop. Monitor inputs for early warning.

% -- FLEX: [CORE] PSI is used in the retraining exercise and incident response
[CORE] The PSI threshold appears again in the P2 severity classification.
IF AHEAD: ``What if your bins are too coarse? Too fine?''
IF SHORT: Show the formula and threshold table, skip the question.
}

\small
\textbf{Population Stability Index (PSI):}
$$\text{PSI} = \sum_{i=1}^{n} (\text{actual}_i - \text{expected}_i) \times \ln\!\left(\frac{\text{actual}_i}{\text{expected}_i}\right)$$

\vspace{-0.1cm}
{\footnotesize
\renewcommand{\arraystretch}{1.0}
\begin{tabular}{@{}lll@{}}
  \toprule
  \textbf{Metric} & \textbf{Threshold} & \textbf{Use Case} \\
  \midrule
  \textbf{PSI}             & $> 0.2$ significant & Categorical / binned \\
  \textbf{KS statistic}    & $> 0.1$ divergent   & Continuous features \\
  \textbf{KL divergence}   & $> 0.1$ significant & Distribution comparison \\
  \bottomrule
\end{tabular}
}

\vspace{0.1cm}
\mlsysinsight{Key:}{Monitor \textbf{input} features --- input shift precedes output shift by the ground-truth feedback loop length.}

\end{frame}

% --- ACTIVE LEARNING 2: Exercise ---
\begin{frame}{Your Turn: Optimal Retraining Interval}
\note{
% -- LINK: From drift metrics to the economics of retraining
PSI tells you drift exists. The $T^*$ formula tells you when to act ---
balancing retraining cost against accuracy loss cost.

% -- NARRATE: Walk through the calculation step by step
``Numerator: 2 times \$500 = \$1,000. Denominator: 10,000 queries/day
times \$0.50 times 0.92 times 0.003 = 13.8. Square root of 1000/13.8
= square root of 72.5 = 8.5 days. Retrain weekly. Daily would be
wasteful. Monthly would violate SLA.''

% -- ENGAGE: Peer Instruction Protocol
Present the problem (30s). Individual calculation (60s).
If 30--70\% get 8.5 days: pair discussion (90s), then re-vote.
[Expected errors: forgetting to multiply all denominator terms]

% -- WARN: Students plug in numbers without understanding the trade-off
Common error: students compute $T^*$ mechanically without seeing the
economic balance. Correct framing: $T^*$ balances retraining expense
vs accuracy decay. If retraining cost drops, $T^*$ shrinks.

% -- FLEX: [CORE] This is the quantitative centerpiece of the chapter
[CORE] The 8.5-day result anchors the Key Takeaways.
IF AHEAD: ``What if lambda doubles?'' ($T^*$ drops to about 6 days.)
IF SHORT: Show the formula and answer directly, skip peer discussion.
}

\footnotesize
\begin{columns}[T]
  \begin{column}{0.55\textwidth}
    {\small\bfseries Calculate $T^*$}

    \vspace{0.1cm}
    A fraud detection model:
    \begin{itemize}\setlength\itemsep{0pt}
      \item $C$ = \$500, \; $Q$ = 10,000/day
      \item $V$ = \$0.50, \; $A_0$ = 0.92
      \item $\lambda$ = 0.003/day
    \end{itemize}

    \vspace{0.05cm}
    $$T^* \approx \sqrt{\frac{2C}{Q \cdot V \cdot A_0 \cdot \lambda}}$$

    {\scriptsize\textcolor{midgray}{(90 seconds --- compare with a neighbor)}}
  \end{column}
  \begin{column}{0.40\textwidth}
    \pause
    \begin{mlsyscard}{datastroke}
    \textbf{Solution:}\\[0.05cm]
    {\scriptsize
    $T^* = \sqrt{\frac{1{,}000}{13.8}} = \sqrt{72.5}$\\[0.05cm]
    $\approx$ \textbf{8.5 days}\\[0.05cm]
    \textcolor{datastroke}{\textbf{Retrain weekly!}}\\
    Daily: wasteful. Monthly: SLA violation.
    }
    \end{mlsyscard}
  \end{column}
\end{columns}

\end{frame}

\begin{frame}{The Uptime Iceberg}
\note{
% -- LINK: From retraining intervals to why monitoring must be layered
The $T^*$ formula assumes you detect drift. This slide shows what happens
when monitoring is shallow.

% -- NARRATE: Walk through the three dashboard panels
Point to each panel: ``Service health: 99.97\% uptime --- green. Model
health: accuracy at 87\%, below the 90\% SLA --- red. Data drift: PSI
approaching 0.2 --- yellow. Traditional monitoring sees only the first
panel. The iceberg is underwater.''

% -- ENGAGE: Challenge the dashboard
Ask: ``Which panel would have caught the ridesharing failure?''
[Expected: Data drift panel --- PSI would have flagged the shift.]

% -- WARN: Students equate uptime with correctness
Common error: students set up infrastructure monitoring and declare the
system monitored. Correct framing: 99.97\% available and perfectly wrong
is worse than a crash --- crashes trigger alerts.

% -- FLEX: [OPTIONAL] Reinforces monitoring theme; not structurally new
[OPTIONAL] The point was made in the opening slide; this visualizes it.
IF SHORT: Show the image and read the alert box. Move on.
}

% --- Layout: FULL-WIDTH IMAGE ---
\centering
\safeimg[width=0.92\textwidth,height=4.8cm,keepaspectratio]{monitoring-dashboard-ml.pdf}

\vspace{0.1cm}
\mlsysalert{Key:}{99.97\% uptime. Model accuracy at 87\%. \textbf{Perfectly available, perfectly wrong.}}

\end{frame}

% =============================================================================
\section{Deployment Strategies}
% =============================================================================

\begin{frame}{Deployment Patterns: Risk vs.\ Cost}
\note{
% -- LINK: From monitoring (detecting problems) to deployment (preventing them)
Monitoring catches drift after deployment. Deployment strategies control
how much damage a bad model can do before monitoring catches it.

% -- NARRATE: Walk through each strategy with cost/risk trade-offs
``Canary: route 5\% of traffic to the new model. Cheapest, but you need
thousands of inferences for statistical confidence. Blue-green: run both,
instant switch. Safe but 2x cost. Shadow: new model runs in parallel but
does not serve users. Safest but 10--20\% overhead.''

% -- ENGAGE: Domain-specific deployment choice
Ask: ``Which strategy for a medical diagnosis model?'' [Expected: shadow,
because the cost of a bad rollout is patient harm.]

% -- WARN: Students default to canary without computing sample size
Common error: students run canary with 100 requests and declare it safe.
Correct framing: ML canary needs statistical significance on model metrics.
With 5\% traffic, reaching significance can take hours or days.

% -- FLEX: [CORE] Deployment strategies connect to rollback and A/B testing
[CORE] The canary pattern recurs in rollback tiers and A/B testing.
IF AHEAD: ``Can you combine strategies? Canary then blue-green?''
IF SHORT: Focus on the diagram and the statistical confidence insight.
}

% --- Layout: FULL-WIDTH IMAGE ---
\centering
\safeimg[width=0.92\textwidth,height=4.8cm,keepaspectratio]{canary-deployment.pdf}

\vspace{0.1cm}
\footnotesize
\mlsysconcept{ML-specific:}{Canary needs \textbf{thousands of inferences} for statistical significance, not just uptime checks.}

\end{frame}

\begin{frame}{Rollback Strategies}
\note{
% -- LINK: From deployment to what happens when deployment fails
Canary, blue-green, and shadow reduce risk. Rollback is the safety net
when a bad model escapes anyway.

% -- NARRATE: Walk through the three tiers with urgency
``Tier 1: Immediate, under 1 minute --- hot standby for crashes. Tier 2:
Rapid, under 15 minutes --- canary metric degradation triggers redeploy.
Tier 3: Delayed, under 4 hours --- business metric decline.'' Then the
fire drill card: ``Untested rollbacks fail at 3 AM. Monthly fire drills
expose stale caches and bad configs before the real incident.''

% -- ENGAGE: Test rollback readiness
Ask: ``When was the last time your team tested a rollback?''
[Expected: most never have. The fire drill rule is the punchline.]

% -- WARN: Students assume rollback = reload old model
Common error: students think rollback is just loading the previous artifact.
Correct framing: rollback involves state (caches, registries), config,
and sometimes data pipeline changes.

% -- FLEX: [OPTIONAL] Table is self-explanatory; fire drill is the insight
[OPTIONAL] The fire drill card is the key takeaway.
IF SHORT: Show the table, read the fire drill card, move on.
}

\footnotesize
\renewcommand{\arraystretch}{1.2}
\begin{tabular}{@{}llll@{}}
  \toprule
  \textbf{Tier} & \textbf{Time} & \textbf{Trigger} & \textbf{Implementation} \\
  \midrule
  \textbf{Immediate} & $<$ 1 min  & Serving errors, crashes & Hot standby, instant switch \\
  \textbf{Rapid}     & $<$ 15 min & Canary metric degradation & Registry-based redeploy \\
  \textbf{Delayed}   & $<$ 4 hrs  & Business metric decline & Full redeploy + state migrate \\
  \bottomrule
\end{tabular}

\vspace{0.2cm}
\begin{mlsyscard}{errorstroke}
\textbf{The Fire Drill Rule}\\
{\footnotesize Rollback procedures that have never been tested \textbf{will fail when needed}.
Monthly fire drills expose configuration gaps, stale caches, and documentation holes
\textbf{before} the 3:00 AM incident.}
\end{mlsyscard}

\end{frame}

\begin{frame}{A/B Testing for ML}
\note{
% -- LINK: From canary (validates stability) to A/B testing (measures impact)
Canary asks ``is the new model safe?'' A/B testing asks ``is it better
for the business?'' Different questions, different methods.

% -- NARRATE: Walk through sample size and decision matrix
``To detect a 2\% lift on a 5\% baseline, you need about 25,000 users
per variant.'' Point to ML challenges: ``Delayed feedback, novelty effects,
interference.'' Point to the decision matrix: ``Improved + guards pass =
ship. Improved + guards fail = investigate.''

% -- ENGAGE: Decision matrix application
Ask: ``Primary metric improved 3\%, but guardrail latency increased 20\%.
Ship or investigate?'' [Expected: investigate --- guardrail failure blocks
shipping even with primary improvement.]

% -- WARN: Students ship on primary metric improvement alone
Common error: students see accuracy improve and ship without checking
guardrails. Correct framing: guardrail metrics are hard constraints;
primary metrics are soft objectives.

% -- FLEX: [OPTIONAL] A/B testing details are supplementary
[OPTIONAL] The decision matrix is the key contribution.
IF SHORT: Skip the sample size equation, cover only the decision matrix.
}

\small
\begin{columns}[T]
  \begin{column}{0.50\textwidth}
    \textbf{Sample Size:}
    $$n = \frac{2(z_{\alpha/2} + z_\beta)^2 \sigma^2}{\delta^2}$$

    {\footnotesize
    For 2\% lift detection, 5\% baseline:\\
    $\approx$ \textbf{25,000 users per variant}
    }

    \vspace{0.1cm}
    \textbf{ML-Specific Challenges:}
    \begin{itemize}\setlength\itemsep{0pt}
      \item {\footnotesize Delayed feedback (days, not seconds)}
      \item {\footnotesize Novelty effects inflate initial results}
      \item {\footnotesize Interference in ranking systems}
    \end{itemize}
  \end{column}
  \begin{column}{0.46\textwidth}
    {\scriptsize
    \renewcommand{\arraystretch}{1.1}
    \begin{tabular}{@{}lll@{}}
      \toprule
      \textbf{Primary} & \textbf{Guards} & \textbf{Decision} \\
      \midrule
      Improved & Pass & Ship \\
      Improved & Fail & Investigate \\
      Neutral  & Pass & Keep current \\
      Degraded & Any  & Do not ship \\
      \bottomrule
    \end{tabular}
    }

    \vspace{0.1cm}
    \begin{mlsyscard}{routingstroke}
    {\footnotesize
    \textbf{Guardrail metrics} must not degrade even if primary metric improves.
    }
    \end{mlsyscard}
  \end{column}
\end{columns}

\end{frame}

% =============================================================================
\section{Incident Response}
% =============================================================================

\begin{frame}{ML Incident Response}
\note{
% -- LINK: From deployment to what happens when things go wrong
Deployment strategies and monitoring reduce risk. Incident response is
the structured process for when failures still occur.

% -- NARRATE: Walk through the debugging tree as a diagnostic algorithm
``Step 1: Is it data? Check PSI and KS. If shifted, data drift. Step 2:
Is it skew? Compare training vs serving distributions. Step 3: Is it
config? Audit recent changes. Step 4: Is it feedback? Check for
self-reinforcing temporal patterns.'' Point to severity table: ``P1:
accuracy drop over 10\%, 30-minute response. P2: localized drift, 4 hours.
P3: latency degradation, 24 hours.''

% -- ENGAGE: Apply the debugging tree
Ask: ``Your model's accuracy dropped 5\% over two weeks. No code changed.
Where do you start?'' [Expected: Step 1 --- check input distributions.]

% -- WARN: Students skip the tree and go straight to retraining
Common error: students detect accuracy loss and immediately retrain.
Correct framing: retraining without diagnosis may entrench the problem.
If the cause is a feedback loop, fresh data includes the corrupted signal.

% -- FLEX: [CORE] The debugging tree is the primary incident response tool
[CORE] This structured approach prevents panic-driven retraining.
IF AHEAD: ``What if steps 1--4 all come back negative?''
IF SHORT: Cover the four steps verbally, skip the severity table.
}

\footnotesize
\textbf{The ML Debugging Tree:}

\vspace{0.05cm}
\begin{enumerate}\setlength\itemsep{0pt}
  \item \textbf{Is it data?} Check input distributions (PSI, KS). Shifted $\to$ data drift.
  \item \textbf{Is it skew?} Compare train vs.\ serve distributions. Diverged $\to$ feature store.
  \item \textbf{Is it config?} Audit recent changes. Found $\to$ config debt.
  \item \textbf{Is it feedback?} Analyze temporal patterns. Self-reinforcing $\to$ loop.
\end{enumerate}

\vspace{0.05cm}
{\scriptsize
\renewcommand{\arraystretch}{1.0}
\begin{tabular}{@{}llll@{}}
  \toprule
  \textbf{Severity} & \textbf{Criteria} & \textbf{Response} & \textbf{Example} \\
  \midrule
  \textbf{P1} & Accuracy drop $>$ 10\% & 30 min & Stale feature data \\
  \textbf{P2} & Localized drift, PSI $>$ 0.3 & 4 hrs & One feature shifted \\
  \textbf{P3} & Latency/throughput degrad. & 24 hrs & Batch misconfigured \\
  \bottomrule
\end{tabular}
}

\end{frame}

% --- ACTIVE LEARNING 3: Discussion ---
\begin{frame}{Discussion: Which Interface Failed?}
\note{
% -- LINK: From the debugging tree to the three-interface framework
The debugging tree diagnoses technical causes. This discussion maps
those causes onto the MLOps interface framework.

% -- NARRATE: Frame the discussion clearly
``A recommendation model dropped from 94\% to 81\% over six months.
No code changed. Infrastructure was green. Which of the three MLOps
interfaces failed?'' Pause. ``No single right answer. All three are
connected. That is the point.''

% -- ENGAGE: Turn-and-talk with structured debrief
Students discuss in pairs for 90 seconds. Cold-call 2--3 pairs.
Expected: Data-Model (feature staleness), Prod-Monitor (no retraining
trigger), or both. Validate all answers and show connections.

% -- FLEX: [CORE] Synthesizes the chapter's three-interface framework
[CORE] The three-interface model is a central organizing principle.
IF SHORT: Do a show-of-hands poll instead of pair discussion.
}

\centering
\vspace{0.6cm}
{\large\bfseries Turn and Talk \textcolor{midgray}{(90 seconds)}}

\vspace{0.5cm}
{\large A recommendation model's accuracy dropped from 94\% to 81\%\\
over six months. No code changed. Infrastructure metrics were green.\\[0.3cm]
\alert{Which of the three MLOps interfaces failed?}}

\vspace{0.5cm}
\begin{columns}[c]
  \begin{column}{0.30\textwidth}\centering\small\textcolor{datastroke}{\textbf{Data-Model}}\\Feature consistency\end{column}
  \begin{column}{0.30\textwidth}\centering\small\textcolor{computestroke}{\textbf{Model-Infra}}\\Environment parity\end{column}
  \begin{column}{0.30\textwidth}\centering\small\textcolor{routingstroke}{\textbf{Prod-Monitor}}\\Retraining cadence\end{column}
\end{columns}

\end{frame}

\begin{frame}{Layered Monitoring Budget}
\note{
% -- LINK: From incident response to prevention through layered monitoring
Incident response is reactive. Layered monitoring makes degradation
visible before it becomes an incident.

% -- NARRATE: Walk through the monitoring layers
``Data freshness: every 15 minutes. Schema: every hour. Feature
distributions: every 4 hours with PSI. Model outputs: every hour.
Business metrics: daily.'' Point to the budget card: ``1 model, 3
variants, 15 metrics, 12 checks/hr = 1,800 data points/hr.''

% -- ENGAGE: Layer selection challenge
Ask: ``Which layer would have caught the session\_length skew?''
[Expected: Feature distribution layer --- PSI on session\_length.]

% -- WARN: Students monitor at the wrong granularity
Common error: students monitor only business metrics (daily) and miss
data drift catchable hourly. Each layer has a different SLA because
failure modes propagate at different speeds.

% -- FLEX: [OPTIONAL] Table is self-explanatory; budget math is the insight
[OPTIONAL] Students can read the table independently.
IF SHORT: Point to the budget card, make the cost-of-monitoring point.
IF AHEAD: ``How would you prioritize with only 3 layers?''
}

\footnotesize
\renewcommand{\arraystretch}{1.15}
\begin{tabular}{@{}llll@{}}
  \toprule
  \textbf{Layer} & \textbf{SLA} & \textbf{Metric Type} & \textbf{Catches} \\
  \midrule
  \textbf{Data freshness}   & 15 min & Staleness, completeness & Stale features \\
  \textbf{Schema}           & 1 hr   & Type/range/null checks  & Pipeline breaks \\
  \textbf{Feature distrib.} & 4 hrs  & PSI, KS, KL divergence  & Data drift \\
  \textbf{Model outputs}    & 1 hr   & Accuracy, confidence     & Model degradation \\
  \textbf{Business metrics}  & 24 hrs & Revenue, engagement     & Downstream impact \\
  \bottomrule
\end{tabular}

\vspace{0.15cm}
\begin{mlsyscard}{computestroke}
{\footnotesize
\textbf{Budget}: 1 model $\times$ 3 variants $\times$ 15 metrics $\times$ 12 checks/hr\\
$=$ \textbf{540 data points/hr} per metric layer\\
Total: $\approx$\textbf{1,800 data points/hr} across all layers
}
\end{mlsyscard}

\end{frame}

% =============================================================================
\section{Maturity Framework}
% =============================================================================

\begin{frame}{MLOps Maturity Levels}
\note{
% -- LINK: From individual practices to organizational maturity
We have covered the components. Maturity is about integrating them.

% -- NARRATE: Walk through the staircase level by level
``Level 0: Jupyter notebooks, manual deployment. Level 1: automated
training, basic CI/CD, model registry. Level 2: full automation with
feature stores, drift detection, canary deployments. Each transition
takes 3--6 months --- not because tools are hard, but integration is.''

% -- ENGAGE: Self-assessment
Ask: ``What maturity level is your team at?'' [Expected: most are
Level 0 or early Level 1.]

% -- WARN: Students think buying tools equals advancing maturity
Common error: students believe installing MLflow moves them to Level 2.
Correct framing: the gap is integration --- connecting the feature store
to the CI/CD pipeline to the drift detector to the deployment strategy.

% -- FLEX: [CORE] The maturity framework organizes all prior content
[CORE] This synthesizes every infrastructure concept into a progression.
IF AHEAD: ``What is the ROI for moving from Level 1 to Level 2?''
IF SHORT: Show the image, state 3--6 months per transition, move on.
}

% --- Layout: FULL-WIDTH IMAGE ---
\centering
\safeimg[width=0.92\textwidth,height=4.8cm,keepaspectratio]{maturity-model.pdf}

\vspace{0.1cm}
\footnotesize
\mlsysconcept{Key insight:}{Each transition takes 3--6 months. The gap is \textbf{integration}, not tools.}

\end{frame}

\begin{frame}{ML Test Score}
\note{
% -- LINK: From qualitative maturity to a concrete scoring rubric
The ML Test Score from Breck et al.\ quantifies maturity: four
categories, specific checks, a numeric score.

% -- NARRATE: Walk through categories and score bands
``Data Tests: distributions verified? Model Tests: accuracy within bounds,
subgroups checked? Infra Tests: training deterministic, rollback works?
Monitor Tests: drift alerts fire?'' Score bands: ``0--5 ad hoc, 15+
production-grade --- very few achieve this.''

% -- ENGAGE: Self-scoring exercise
Ask: ``Where would your team score?'' [Expected: most below 5.]

% -- WARN: Students overestimate their maturity
Common error: having a model registry does not mean scoring 10+.
The test score requires verified, automated checks actually running.

% -- FLEX: [OPTIONAL] Rubric reinforces maturity; not structurally new
[OPTIONAL] Students can study the rubric independently.
IF SHORT: Show table and score bands, skip self-scoring.
IF AHEAD: ``Pick one category and design a test suite for it.''
}

\scriptsize
\renewcommand{\arraystretch}{1.05}
\begin{tabular}{@{}llp{5cm}@{}}
  \toprule
  \textbf{Category} & \textbf{Tests} & \textbf{Example Checks} \\
  \midrule
  \textbf{Data Tests}    & Feature expectations, schema validation  & Distributions verified; no training-serving skew \\
  \textbf{Model Tests}   & Staleness, fairness, performance bounds  & Accuracy within bounds; no subgroup degradation \\
  \textbf{Infra Tests}   & Reproducibility, pipeline reliability    & Training is deterministic; rollback works \\
  \textbf{Monitor Tests} & Alerts, dashboards, incident response    & Drift alerts fire; on-call rotation documented \\
  \bottomrule
\end{tabular}

\vspace{0.15cm}
\small
\begin{columns}[T]
  \begin{column}{0.22\textwidth}
    \centering
    \textcolor{errorstroke}{\textbf{0--5}}\\{\scriptsize Ad hoc}
  \end{column}
  \begin{column}{0.22\textwidth}
    \centering
    \textcolor{routingstroke}{\textbf{5--10}}\\{\scriptsize Developing}
  \end{column}
  \begin{column}{0.22\textwidth}
    \centering
    \textcolor{computestroke}{\textbf{10--15}}\\{\scriptsize Mature}
  \end{column}
  \begin{column}{0.22\textwidth}
    \centering
    \textcolor{datastroke}{\textbf{15+}}\\{\scriptsize Production-grade}
  \end{column}
\end{columns}
\mlsyscite{Breck et al., ICSE 2020}

\end{frame}

% =============================================================================
\section{Wrap-Up}
% =============================================================================

\begin{frame}{Fallacies}
\note{
% -- NARRATE: Read each fallacy with its quantitative rebuttal
``Fallacy 1: MLOps is just DevOps for ML. Reality: ML pipelines take
15--45 min vs 2--5 for traditional CI/CD.'' ``Fallacy 2: Automated
retraining ensures performance. Reality: automation propagates corrupted
data.'' ``Fallacy 3: Consistency is automatic. Reality: 5--15\% accuracy
loss even with established pipelines.''

% -- FLEX: [CORE] Fallacies crystallize the chapter's key warnings
[CORE] Each fallacy maps to a concept from the lecture.
IF SHORT: Cover fallacies 1 and 3, skip fallacy 2.
}

\small
\textbf{Fallacy:} \textit{MLOps is just applying DevOps to ML models.}\\
{\footnotesize ML pipelines take 15--45 min vs.\ 2--5 min for traditional software. They require data validation, model evaluation, and drift detection stages absent from standard CI/CD.}

\vspace{0.15cm}
\textbf{Fallacy:} \textit{Automated retraining ensures optimal performance without oversight.}\\
{\footnotesize Automated pipelines propagate corrupted training data. A news model retrained on weekend data showed lower engagement because user behavior differs sharply weekday vs.\ weekend.}

\vspace{0.15cm}
\textbf{Fallacy:} \textit{Training-serving consistency is automatic once pipelines are established.}\\
{\footnotesize Skew causes 5--15\% accuracy degradation even with established pipelines. An e-commerce model lost 12\% accuracy for 6 weeks from a \texttt{session\_length} computation mismatch.}

\end{frame}

\begin{frame}{Pitfalls}
\note{
% -- NARRATE: Read each pitfall with quantitative evidence
``Pitfall 1: Treating deployment as one-time. Fraud model PSI went from
0.18 to 0.31 in 3 months, accuracy 94\% to 87\%. MTTR: 5--14 days
without automation, 4--24 hours with it.'' ``Pitfall 2: Monitoring outputs
without inputs. 18-hour stale embeddings caused 9\% degradation for 3 days
before alerts fired.'' ``Pitfall 3: Infrastructure without organizational
alignment. Siloed teams deploy quarterly; unified teams deploy weekly.''

% -- FLEX: [CORE] Pitfalls provide concrete operational warnings
[CORE] Each pitfall reinforces a chapter theme.
IF SHORT: Cover pitfalls 1 and 2, skip pitfall 3.
}

\small
\textbf{Pitfall:} \textit{Treating deployment as a one-time event.}\\
{\footnotesize A fraud model with PSI = 0.18 at deployment reached PSI = 0.31 within 3 months, dropping from 94\% to 87\% accuracy. Without automated retraining, MTTR averages 5--14 days; with automation, 4--24 hours.}

\vspace{0.15cm}
\textbf{Pitfall:} \textit{Monitoring outputs without validating inputs.}\\
{\footnotesize A recommendation system tracked CTR but ignored feature staleness. User embeddings were 18 hrs stale, causing 9\% degradation for 3 days before accuracy alerts triggered.}

\vspace{0.15cm}
\textbf{Pitfall:} \textit{Investing in infrastructure while neglecting organizational alignment.}\\
{\footnotesize A retail company with feature stores and registries still deployed quarterly because data scientists and engineers operated in isolation. Unified teams deploy weekly.}

\end{frame}


\begin{frame}{Muddiest Point}
\note{
% -- NARRATE: Cue the one-minute paper
``Take 60 seconds. Write the muddiest point --- the concept that felt
most confusing. Hand it in on your way out.''

% -- FLEX: [CORE] Drives next lecture's opening clarification
[CORE] Collect responses. Open Ch15 with the top 2--3 muddiest points.
IF SHORT: Allow 30 seconds instead of 60.
}

\centering
\Large\bfseries One-minute paper\\[0.5cm]
\normalsize Write down the \textbf{muddiest point} from today's lecture.\\[0.2cm]
{\small What concept was most confusing or unclear?}\\[0.5cm]
{\footnotesize\textcolor{midgray}{(Hand in on your way out --- or submit digitally)}}
\end{frame}

% --- RETRIEVAL PRACTICE ---
\begin{frame}{What Were the Key Ideas?}
\note{
% -- NARRATE: Cue retrieval practice with silence
``Close your notes. 90 seconds. Write the 4 most important concepts.
No peeking.'' Walk around the room. Do NOT show Key Takeaways yet.

% -- FLEX: [CORE] Retrieval practice is the highest-impact activity
[CORE] Forces active recall before the reveal.
IF SHORT: Reduce to 60 seconds.
}

\centering
\vspace{1.5cm}
{\Large\bfseries Close your notes.}

\vspace{0.8cm}
{\large Write down the \textbf{4 most important concepts} from today.}

\vspace{0.8cm}
{\normalsize\textcolor{midgray}{90 seconds --- no peeking.}}

\end{frame}

\begin{frame}{Key Takeaways}
\note{
% -- LINK: From retrieval practice to the official summary
Students just wrote their own summaries. Now compare with the canonical list.

% -- NARRATE: Walk through each bullet with its quantitative anchor
``Silent failures: the Degradation Equation. Skew: \$365K/year. Retraining:
$T^* = 8.5$ days. Rollout: three tiers. Monitor inputs. Five principles.
Maturity: 3--6 months per transition.''

% -- FLEX: [CORE] Key Takeaways are the exam-relevant summary
[CORE] Students photograph this slide. Read each bullet clearly.
IF SHORT: Read bullets without elaboration.
}

\scriptsize
\begin{itemize}\setlength\itemsep{0pt}
  \item \textbf{ML systems fail silently}: Degradation Equation $\text{Acc}(t) \approx \text{Acc}_0 - \lambda \cdot D(P_t \| P_0)$ quantifies decay.
  \item \textbf{Training-serving skew}: Feature stores eliminate 5--15\% accuracy loss by computing features once.
  \item \textbf{Retraining is an optimization}: $T^* \approx \sqrt{2C/(QVA_0\lambda)}$ --- from intuition to economics.
  \item \textbf{Graduated rollout}: Canary, blue-green, shadow with tiered rollback ($<$1 min, $<$15 min, $<$4 hrs).
  \item \textbf{Monitor inputs, not just outputs}: Input drift precedes output degradation.
  \item \textbf{Five Principles}: Reproducibility, Separation, Consistency, Observable Degradation, Cost-Aware Automation.
  \item \textbf{Maturity is integration, not tools}: Each level transition takes 3--6 months.
\end{itemize}

\end{frame}

\begin{frame}{References}
\note{
% -- NARRATE: Highlight the two must-read papers
``Sculley et al.\ 2015 on ML technical debt. Breck et al.\ 2020 for the
ML Test Score rubric.''

% -- FLEX: [OPTIONAL] Reference slides are for student self-study
[OPTIONAL] Do not read every reference aloud.
IF SHORT: Name two key papers and move to Next Lecture.
}

\small
\mlsysref{Sculley+15}{Sculley et al. ``Hidden Technical Debt in ML Systems.'' NeurIPS 2015.}
\mlsysref{Breck+20}{Breck et al. ``The ML Test Score: A Rubric for Production Readiness.'' ICSE 2020.}
\mlsysref{Amershi+19}{Amershi et al. ``Software Engineering for ML: A Case Study.'' ICSE 2019.}
\mlsysref{Hermann+17}{Hermann \& Del Balso. ``Meet Michelangelo: Uber's ML Platform.'' 2017.}
\mlsysref{Paleyes+22}{Paleyes et al. ``Challenges in Deploying ML.'' ACM Comp.\ Surveys, 2022.}

\end{frame}

\begin{frame}{Next Lecture: Responsible Engineering}
\note{
% -- LINK: From operational correctness to ethical correctness
This chapter ensured systems stay correct. The next asks: correct for whom?

% -- NARRATE: Deliver the forward hook with the three pillars
``We built a system that is efficient, scalable, and reliable. Can it
still cause harm?'' Point to the three pillars: ``Bias, Privacy, Trust.''
Close: ``The final constraint is aligning optimization with human values.''

% -- ENGAGE: Provocative closing question
Ask: ``Can a system be perfectly correct and still unfair?''
[Expected: yes --- COMPAS was calibrated but had disparate error rates.]

% -- FLEX: [CORE] The forward hook creates anticipation for Ch15
[CORE] End on this question. Do not answer it --- let it simmer.
IF SHORT: Read the slide text and close.
}

\centering
\vspace{0.5cm}
{\large We built a system that is \textbf{efficient}, \textbf{scalable}, and \textbf{reliable}.}

\vspace{0.5cm}
{\normalsize Can it still cause harm?}

\vspace{0.4cm}
\begin{columns}[c]
  \begin{column}{0.30\textwidth}
    \centering
    \textcolor{datastroke}{\Large\bfseries Bias}\\[0.1cm]
    {\small Amplifying inequality}
  \end{column}
  \begin{column}{0.30\textwidth}
    \centering
    \textcolor{errorstroke}{\Large\bfseries Privacy}\\[0.1cm]
    {\small Leaking data}
  \end{column}
  \begin{column}{0.30\textwidth}
    \centering
    \textcolor{routingstroke}{\Large\bfseries Trust}\\[0.1cm]
    {\small Explainability gap}
  \end{column}
\end{columns}

\vspace{0.4cm}
{\normalsize\textbf{The final constraint: aligning optimization with human values.}}

\end{frame}


% =============================================================================
% BACKUP SLIDES
% =============================================================================
\appendix

\begin{frame}{Backup: PSI Threshold Decision Guide}
\note{
% -- NARRATE: Walk through each PSI band with the action
PSI below 0.1: normal variation. PSI 0.1--0.2: investigate which features
shifted. PSI above 0.2: trigger retraining. PSI above 0.5: likely a
pipeline failure, not organic drift.

% -- FLEX: [OPTIONAL] Use if students struggled with PSI thresholds
[OPTIONAL] Pull if the main PSI slide generated many questions.
}
\small
Use if students need a practical drift detection reference.\\small PSI $< 0.1$: No action. Normal variation.\\ PSI $0.1$--$0.2$: Investigate. Check if specific features shifted.\\ PSI $> 0.2$: Significant drift. Trigger retraining pipeline.\\ PSI $> 0.5$: Critical. Likely data pipeline failure, not organic drift.\\ Monitor input features, not just output accuracy.
\end{frame}

\begin{frame}{Backup: Retraining Interval Optimization}
\note{
% -- NARRATE: Explain $T^*$ sensitivity to each parameter
High C: retrain less often. High Q: retrain more often. High lambda:
retrain more often. The formula balances retraining expense against
accuracy loss cost.

% -- FLEX: [OPTIONAL] Use if students want more practice with $T^*$
[OPTIONAL] Pull for additional worked examples or if the exercise
generated confusion about parameter sensitivity.
}
\small
Use if students want more practice with the $T^*$ formula.\\small $$T^* \approx \sqrt{\frac{2C}{Q \cdot V \cdot A_0 \cdot \lambda}}$$\\ High retraining cost ($C$ large): retrain less often.\\ High query volume ($Q$ large): retrain more often.\\ Fast drift ($\lambda$ large): retrain more often.\\ The formula balances retraining cost against accuracy loss cost.
\end{frame}

\end{document}
